\documentclass[11pt,a4paper]{article}
\usepackage{fontspec}\usepackage{polyglossia}
\setdefaultlanguage{russian}\setotherlanguage{english}
\newcommand{\optserif}{Path=fonts/,Extension=.ttf,UprightFont=*-Regular,BoldFont=*-Bold,ItalicFont=*-Italic,BoldItalicFont=*-BoldItalic,Ligatures=TeX}
\IfFileExists{fonts/LiberationSerif-Regular.ttf}{%
  \setmainfont{LiberationSerif}[\optserif]\newfontfamily\cyrillicfont{LiberationSerif}[\optserif]}{%
  \IfFontExistsTF{Liberation Serif}{\setmainfont{Liberation Serif}\newfontfamily\cyrillicfont{Liberation Serif}}{%
    \setmainfont{DejaVu Serif}\newfontfamily\cyrillicfont{DejaVu Serif}}}
\usepackage[a4paper,margin=2.0cm]{geometry}
\usepackage{enumitem,amsmath,amssymb,booktabs,titlesec,amsthm}
\usepackage[hidelinks]{hyperref}
\setlist[enumerate]{leftmargin=1.3em,itemsep=1pt,topsep=2pt}
\setlist[itemize]{leftmargin=1.1em,itemsep=1pt,topsep=2pt}
\setlength{\parskip}{2pt}\emergencystretch=2em
\titleformat{\section}{\large\bfseries}{\thesection.}{0.4em}{}
\titlespacing*{\section}{0pt}{8pt plus 1pt}{3pt}
\newtheorem{prop}{Предложение}
\begin{document}
\begin{center}
{\large\bfseries Отображение окружности переключателя сна: вывод в замкнутой форме}\\[2pt]
{\small Приложение к отчёту о программе опровержения. 30 октября 2026 г.}
\end{center}
\vspace{-2pt}

\noindent\textbf{Что здесь показано.} Слой сна и бодрствования модели сводится \emph{точно} к реле
с гистерезисом, и его отображение последования выписывается в замкнутой форме с точностью до двух
трансцендентных уравнений. Отсюда: складки быстрой подсистемы в явном параметрическом виде
($1{,}4503$ и $2{,}4635$~мВ против $1{,}450/2{,}463$ в модели); свободный период в замкнутой форме;
край языка~$1{:}1$ как решение системы из трёх уравнений, без всякого перебора. И отсюда же следует
ответ на вопрос об инвариантных торах: их здесь \emph{нет}, и это доказывается неравенством на
показатель Ляпунова, а не перебором. Попутно оказывается, что рамка языков Арнольда к этой системе
неприменима: отображение кусочно-гладкое и окружность не накрывает.

\section{Редукция к реле с гистерезисом}
Быстрая пара напряжений ядер ($\tau\approx10$~с) при временах сомногенов (18{,}2 и 4{,}2~ч)
исключается адиабатически. Её стационар допускает \emph{параметрическое} решение: взяв $V_m$ за
параметр,
\begin{equation}\label{eq:par}
Q(V_v)=\frac{A_m-V_m}{|\nu_{mv}|},\qquad V_v=Q^{-1}\!\Bigl(\frac{A_m-V_m}{|\nu_{mv}|}\Bigr),\qquad
D_v(V_m)=V_v-\nu_{vm}Q(V_m),
\end{equation}
где $Q^{-1}(q)=\theta+\sigma_Q\ln\bigl(q/(Q_{\max}-q)\bigr)$. Складки суть экстремумы
$D_v(V_m)$. Численно по~\eqref{eq:par}: $D_-=1{,}4503$~мВ (ветвь сна, $Q(V_m)=0{,}75$) и
$D_+=2{,}4635$~мВ (ветвь бодрствования, $Q(V_m)=3{,}45$) --- против $1{,}450$ и $2{,}463$,
заявленных в модели как сверка с публикацией. Совпадение до четвёртой значащей цифры получено
\emph{без интегрирования}.

В пределе $\tau\to0$ система есть реле: сомногены подчиняются линейным уравнениям с замкнутыми
решениями, а переключение происходит на двух \emph{подвижных порогах}
\begin{equation}\label{eq:thr}
g_\pm(t)=\frac{D_\pm-B_0+\kappa-\nu_{vc}\cos\bigl(2\pi(t-\alpha)/24\bigr)}{A_G},
\qquad \kappa=k\,D .
\end{equation}

\begin{prop}[Отображение последования]\label{prop:map}
Пусть $t_n$ --- момент засыпания. Тогда $t_{n+1}=f(t_n)$ задано неявно системой
\begin{equation}\label{eq:map}
g_+(t_n)\,e^{-(t_w-t_n)/\tau_s}=g_-(t_w),\qquad
1-\bigl(1-g_-(t_w)\bigr)e^{-(t_{n+1}-t_w)/\tau_w}=g_+(t_{n+1}),
\end{equation}
где $t_w$ --- момент пробуждения. Оба уравнения имеют вид «экспонента = синусоида». Отображение
гладко там, где пересечение трансверсально, и теряет гладкость там, где оно касательное; последнее
условие и есть критическая линия.
\end{prop}

\section{Свободный период и поправка на задержанную бифуркацию}
При $\nu_{vc}=0$ пороги постоянны, и~\eqref{eq:map} решается явно:
\begin{equation}\label{eq:per}
T_s=\tau_s\ln\frac{g_+}{g_-},\qquad
T_w=\tau_w\ln\frac{1-g_-}{1-g_+},\qquad T_0=T_s+T_w .
\end{equation}
Подстановка даёт $T_s=1{,}337$~ч, $T_w=6{,}683$~ч, $T_0=8{,}020$~ч. Прямой счёт при нулевой
амплитуде даёт $8{,}878$~ч --- расхождение 10~\%. Причина найдена и она не ошибка вывода:
переключение происходит \emph{не на складке}. Засыпание наступает при $D_v=2{,}5026$ против складки
$2{,}4635$, пробуждение при $1{,}3794$ против $1{,}4503$: система проскакивает складку, поскольку у
седло-узла время релаксации расходится и медленный проход завершается позже точки бифуркации.
Проскок измерен на двух порядках по быстрому времени и обращается в нуль с показателем $0{,}64$ и
$0{,}59$ при теоретическом $2/3$ для медленного прохода через седло-узел. Значит~\eqref{eq:per}
есть \emph{точный предел} $\tau\to0$, а модель стоит в 10~\% от него; при использовании
аналитических порогов это смещение обязано быть названо, а не скрыто.

\section{Край языка 1:1 без перебора}
Орбита $1{:}1$ есть неподвижная точка $f$ со сдвигом на сутки. Край языка --- её седло-узел:
\begin{equation}\label{eq:edge}
f(t^*)=t^*+24,\qquad f'(t^*)=1,
\end{equation}
что вместе с~\eqref{eq:map} даёт три уравнения на $(t^*,t_w,\kappa)$ и определяет
$\kappa_{\mathrm{edge}}$ с машинной точностью. Численная биссекция по признаку захвата дала
$4{,}7651$~мВ, независимый перебор числа вращения (правило~Ф20) --- $4{,}77$~мВ; расхождение
$0{,}005$~мВ. Границы языков $p{:}q$ получаются так же: периодическая орбита периода $q$ плюс
условие на мультипликатор.

\section{Инвариантных торов нет: показатель отделён от нуля}
\textbf{Отображение определено не всюду.} Сомногены растут к единице снизу, поэтому засыпание
возможно лишь там, где порог $g_+(t)<1$. Доля таких фаз:
\begin{center}\small
\begin{tabular}{lccccc}
\toprule
$\kappa$, мВ & 0 & 2 & 4{,}80 & 5{,}10 & 5{,}18\\
доля суток, где засыпание возможно & 0{,}789 & 0{,}531 & 0{,}166 & 0{,}077 & 0{,}016\\
\bottomrule
\end{tabular}
\end{center}
При $\kappa\approx5{,}18$ она обращается в нуль --- так исчезает сон. При $\kappa=4{,}80$ область
определения есть дуга $[10{,}014;\,13{,}986]$~ч длиной $3{,}97$, а образ --- дуга
$[10{,}846;\,11{,}839]$ длиной $0{,}993$, и образ лежит \emph{внутри} области.

Отсюда сразу следует, что классическая рамка языков Арнольда к этой системе \textbf{неприменима}.
У отображения окружности степени один образ обязан накрывать всю окружность; здесь он занимает
4~\% её. Отображение кусочно-гладкое, с разрывами: при $\kappa=4{,}88$ скачок составляет
$22{,}47$~ч при медианном шаге $10^{-4}$~ч, то есть почти ровно сутки. Это и есть пропущенная ночь.
\emph{Лестница захвата порождается сменой ветви, а не захватом гладкого отображения окружности.}

\textbf{Правильная величина --- показатель вдоль орбиты.} Максимум $|f'|$ по сетке для этого не
годится: у касательных пересечений производная имеет корневую особенность, и сетка ловит точки, куда
аттрактор не заходит. Итерация аналитического отображения~\eqref{eq:map} с отбрасыванием переходного
участка даёт:
\begin{center}\small
\begin{tabular}{lcccccc}
\toprule
$\kappa$, мВ & 4{,}78 & 4{,}82 & 4{,}88 & 4{,}94 & 5{,}00 & 5{,}10\\
$\lambda$ вдоль орбиты & $-3{,}17$ & $-3{,}01$ & $-3{,}24$ & $-3{,}34$ & $-3{,}54$ & $-2{,}53$\\
число вращения $\rho$, сут & 1{,}0000 & 1{,}5000 & 2{,}0000 & 2{,}0000 & 2{,}0000 & 2{,}0000\\
\bottomrule
\end{tabular}
\end{center}
По всему окну $\lambda$ лежит между $-1{,}62$ и $-3{,}54$. Отсюда два вывода, и оба --- неравенства,
а не наблюдения. \emph{Хаос исключён}: он требует $\lambda>0$. \emph{Гладкий инвариантный тор
исключён}: квазипериодический аттрактор требует $\lambda=0$ ровно, а здесь расстояние до нуля не
меньше $1{,}6$. Это не «не нашли», а «далеко». Аттрактор при каждом проверенном $\kappa$ есть
сильно притягивающая периодическая орбита.

Числа вращения $1$, $3/2$, $2$ воспроизводят лестницу правила~Ф20 ($1{:}1\to5{:}4\to3{:}2\to2{:}1$;
узкая ступень $5{:}4$ в двенадцать точек не попала). Это третье независимое подтверждение редукции
--- после складок и после края языка.

\textbf{Чего нельзя было сделать сеткой, и это стоит записать отдельно.} В первой редакции
правила~Ф30 часть~(ж) гласила: если доля захвата на сетке меньше единицы, то торы существуют на
множестве параметров положительной меры. Вывод неверен. Дополнение к захваченному множеству всюду
плотно, поэтому конечная сетка найдёт незахваченные узлы и при положительной мере дополнения, и при
нулевой, а промахнётся мимо них в обоих случаях с вероятностью, зависящей только от разрешения:
мощности у такого критерия нет ни в одну сторону. Честное содержание результата «240 узлов из 240
захвачены» --- только граница снизу на меру захвата, $0{,}988$ при доверии 95~\%. Вопрос же решается
не мерой, а показателем, и решается в одну строку.

Путь к правильной величине занял два неверных объекта: сперва доля захвата по сетке, затем максимум
$|f'|$ по области. Оба измеряли не то. Содержательный урок общий: на сечении Пуанкаре величина,
усреднённая \emph{вдоль орбиты}, и величина, взятая по \emph{всей области}, отвечают на разные
вопросы, и вторая к аттрактору отношения может не иметь.

\textbf{Мера в фазовом пространстве --- вопрос другой, и здесь он не возникает.} Всё сказанное выше
относится к пространству \emph{параметров}. Вопрос о мере самого аттрактора --- о толстом
аттракторе положительной меры Лебега, когда его размерность совпадает с размерностью системы, ---
принадлежит другому и существенно более трудному кругу задач (примеры для ступенчатых косых
произведений строил Ильяшенко). Наша система стоит на противоположном полюсе, и это измерено.
Быстрые времена переключателя (10~с) против медленных (часы) дают огромно отрицательную дивергенцию,
и две быстрые размерности схлопываются сразу: контроль~Ф30(а) дал одномерность отображения
последования с точностью $4{,}5\cdot10^{-13}$ --- восемь траекторий с разных начальных сомногенов
легли на одну кривую до машинного нуля. Само одномерное отображение притягивает к периодической
орбите с $\lambda\le-1{,}6$, то есть аттрактор в сечении есть конечное множество точек, а в
расширенном фазовом пространстве --- замкнутая кривая. Это следствие устройства: чтобы приблизиться
к толщине, дивергенция должна где-то перестать быть отрицательной, то есть быстрые времена обязаны
стать сравнимыми с медленными, --- но именно на их разделении держится вся редукция к реле, а с ней
складки и отображение окружности.

\section{Что из этого следует для счёта}
Построение отображения по~\eqref{eq:map} требует двух скалярных корней на цикл вместо 1600 шагов на
сутки, а край языка по~\eqref{eq:edge} --- решения системы из трёх уравнений вместо биссекции.
Прогон Ф30 занял 1686~с; аналитический путь даёт те же числа за секунды, с точностью до названного
10~\%-го смещения порогов. Отдельно это касается части~(е), которая в прогоне осталась
непроверяемой: показатель вдоль орбиты считается двенадцатью итерациями отображения и отвечает на её
вопрос прямо, тогда как сетка из 240 узлов на него ответить не могла в принципе --- она искала
непериодические узлы там, где аттрактор периодичен при каждом параметре.

\smallskip\noindent\textbf{Открытые места.} (i)~Показатель проскока сошёлся к $0{,}59$--$0{,}64$ при
теоретическом $2/3$; асимптотика не достигнута, и либо нужен ещё порядок по быстрому времени, либо
поправка содержит логарифм. (ii)~Показатель вдоль орбиты посчитан в двенадцати точках по $\kappa$;
формально это не исключает исключительных значений, на которых мог бы жить канторов минимальный
аттрактор типа Данжуа. Довод против --- структурный, а не статистический: такому аттрактору тоже
нужен показатель у нуля, а он нигде не поднимается выше $-1{,}6$. (iii)~Разделение ветвей
отображения сделано порогом на величину скачка и надёжным не является; число ветвей в тексте не
приводится намеренно. (iv)~Всё изложенное относится к детерминированному переключателю с замороженным
повреждением: при живом повреждении и шуме нагрузки речь о аттракторах вообще не идёт.
\end{document}
